\begin{table}[h!]
\caption{Literature survey (n=21). Model-free baseline cells name the untrained comparator and cite its original paper, not the surveyed one. The first six rows are the pilot sample.}
\label{tab:literature_survey}
\footnotesize
\resizebox{\textwidth}{!}{%
\begin{tabular}{llcll}
\hline
Paper & Task & Held-out edges & Negative sampling & Model-free \\
 & & removed from graph? & & baseline \\
\hline
MGCNSS \cite{tian2024mgcnss} & miRNA-disease & no & distance-based selection (paper's own contribution) & no \\
NIMGSA \cite{jin2022nimgsa} & miRNA-disease & unclear$^{\ddagger}$ & all unlabeled pairs (implicit: full-matrix loss) & SPM~\cite{lu2015spm} \\
HybridGNN \cite{ahmad2026hybridgnn} & miRNA-disease & yes$^{\ddagger}$ & uniform random from unknown pairs & no \\
HGDTI \cite{yu2022hgdti} & drug-target & unclear$^{\ddagger}$ & reliable score-filtered (non-uniform) & no \\
NGCN \cite{cao2024ngcn} & drug-target & yes$^{\ddagger}$ & uniform random from unknown pairs, 1:10 & HNM~\cite{wang2014hnm} \\
kmerPMTF \cite{zhang2024kmerpmtf} & plant miRNA-target & unclear$^{\ddagger}$ & all unlabeled pairs, count-matched & no \\
\hline
Orro \cite{orro2026mirnadisease} & miRNA-disease & unclear$^{\S}$ & uniform random from unannotated pairs & TCRWMDA~\cite{yu2020tcrwmda} \\
CoupleMDA \cite{li2025couplemda} & miRNA-disease & yes$^{\ddagger}$ & uniform random, 1:1 & no \\
GONNMDA \cite{zeng2025gonnmda} & miRNA-disease & unclear$^{\P}$ & uniform random, 1:1 & no \\
DiGAMN \cite{liu2024digamn} & miRNA-disease & yes & uniform random, 1:1 / 1:5 / 1:10 & no \\
DGNMDA \cite{lu2024dgnmda} & miRNA-disease & unclear & uniform random from 0-labeled pairs, 1:1 & no \\
HLGNN-MDA \cite{yu2022hlgnnmda} & miRNA-disease & yes$^{\ddagger}$ & uniform random, 1:1 & BNPMDA~\cite{chen2018bnpmda} \\
MEAHNE \cite{huang2022meahne} & miRNA-disease & unclear$^{\ddagger}$ & uniform random, 1:1 & no \\
CKSNP-GNN \cite{li2022cksnpgnn} & miRNA-disease & unclear$^{\ddagger}$ & uniform random, 1:1 & LLCMDA~\cite{qu2018llcmda} \\
HMCDA \cite{liang2023hmcda} & circRNA-disease & unclear$^{\ddagger}$ & uniform random, 5:1 & no \\
HiGLDP \cite{wang2026higldp} & lncRNA-disease & unclear$^{\ddagger}$ & not described (1:1 ratio only) & no \\
GPS-DTI \cite{xiong2025gpsdti} & drug-target & n/a$^{\dagger}$ & not described (balanced; 1--20$\times$ in the COVID-19 arm only) & no \\
SaeGraphDTI \cite{zhang2025saegraphdti} & drug-target & yes & not described (balanced sets inherited from prior work) & no \\
DTI-MHAPR \cite{yang2025dtimhapr} & drug-target & unclear & uniform random, 1:1 & named (FSI), no result \\
iNGNN-DTI \cite{sun2024ingnndti} & drug-target & n/a$^{\dagger}$ & uniform random, 1:1 & no \\
ModulePred \cite{jia2024modulepred} & gene-disease & unclear & uniform random, 50:1 & RWR~\cite{tong2006rwr}, RWRH~\cite{li2010rwrh} \\
\hline
\end{tabular}
}

\medskip\par $^{\dagger}$iNGNN-DTI~\cite{sun2024ingnndti} and GPS-DTI~\cite{xiong2025gpsdti} encode drug and protein separately and never
propagate over a graph containing the predicted associations, so the held-out-edges
question does not apply to them.
$^{\ddagger}$Every row is classified from the paper's text. NIMGSA's released code shows what its text
does not: its folds remove the held-out miRNAs' rows from the model's input, but it scores the whole
matrix, training rows included, over a similarity graph built from all associations
(Section~\ref{sec:generalization}). HybridGNN's text describes a held-out edge split, but its released
code builds the message-passing graph and its features from the full matrix, test edges included
(Section~\ref{sec:generalization}). HGDTI's released code propagates at test time over an adjacency built from
the test labels (\ref{app:nimgsa}). NGCN's released code confirms its row: its model reads no
drug--target edge (\ref{app:nimgsa}). kmerPMTF's released code evaluates over a graph that holds the
held-out positives as edges. CoupleMDA's released code confirms its row, except that 42 of its 1,523 test
positives are also training edges; its negatives keep each positive's disease and draw the miRNA at random, not
uniform pairs as its text says (\ref{app:nimgsa}). HLGNN-MDA's released code confirms its row: it masks the test
positives in both directions and reports the last epoch, though its script runs one 90/10 split, not the
tenfold cross-validation its text reports. MEAHNE's released code cannot settle its row: it reads metapath
neighbour lists and splits that it does not release, and the adjacency it does release holds every association. CKSNP-GNN's released code propagates at test time over every sampled pair,
test positives and negatives alike, so the graph does not separate them, but its disease features are 96\% a
Gaussian interaction-profile kernel of the full association matrix. HMCDA's released code removes test
associations from its metapath graph, but 168 duplicate rows in its association file leave 56 of its 432 test
positives (13\%) as training edges, and it trains on one negative per positive, not the five its text reports. HiGLDP's released code keeps test pairs out of its pair graph, but one of the four lncRNA similarity networks behind its features is a Gaussian interaction-profile kernel of the full association matrix, test associations included; its negatives are a fixed uniform draw from unknown pairs, and it reports each fold's best test-epoch AUROC.
$^{\S}$Orro~\cite{orro2026mirnadisease} holds out whole miRNAs, the strongest split in the survey, but removes
``all associations involving training miRNAs \ldots{} from the validation graph'' and builds its
pattern--disease edges ``symmetrically for the validation graph'', so the graph it validates on may
hold the held-out miRNAs' own associations; the text does not settle which, and no code is released.
$^{\P}$GONNMDA's~\cite{zeng2025gonnmda} code link points to an empty repository, and its features include miRNA
functional similarity and Gaussian interaction-profile kernels computed from the association matrix, with no
statement that they are recomputed per fold.
\end{table}
